%% GENERATED by python/paper/build.py from results/ -- do not edit by hand.
%% Rebuild:  .venv\Scripts\python.exe python\paper\build.py --only ArgentinaCalibration_vectorX
%% Source:   results/paper/calibrationSummary.csv
\begin{table}[!htb]
\centering
\begin{threeparttable}
\caption{Calibration, Argentina (vector $X_i$ calibration)}
\label{table:Arg:Calib_vectorX}
\begin{tabular}{lll}
\toprule
\textbf{Parameter} & \textbf{Value} & \textbf{Target} \\
\midrule 
$\epsilon$ (\textit{pre}-reform) & $0.21$ & Basic pension, 70\% of the minimum, discounted \\[1ex] % row: epsilon
$\theta$ & $0.84$ & Replacement rate dispersion \\[1ex] % row: theta
$\alpha$ & $0.35$ & Factor income shares \\[1ex] % row: alpha
$\nu_t$ & $[1.65, 0.96]$ & 30-year gross population growth rates \\[1ex] % row: nu
$\xi$ & $0.30$ & Elasticity of labor supply \\[1ex] % row: xi
$\beta$ & $0.65$ & Capital--output ratio of $3.23$ \\[1ex] % row: beta
$X_i$ & $[0.74, 0.77, 0.88, 1.24]$ & Relative working hours \\[1ex] % row: X
$\eta_i$ & $[0.51, 0.72, 0.95, 1.69]$ & Income distribution \\[1ex] % row: eta
$\gamma_0$ & $0.47$ & Elderly without a pension, 32\% of all households \\[1ex] % row: gamma0
$\omega$ & $1.45$ & Social security tax of $10.9\%$ \\[1ex] % row: omega
$\eta_0$ & $0.288$ & Informal relative income, eq.\ \eqref{eq:calibration_eta} \\[1ex] % row: eta0
$X_0$ & $0.375$ & Informal relative working hours \\[1ex] % row: X0
\bottomrule
\end{tabular}
\begin{tablenotes}[flushleft]
\footnotesize
\item[] \textit{Note:} Our default specification relies on $\rho=1$. Vector-$X_i$ calibration: $X_i$ is identified from relative hours, which are data here, and the level of $\bar h$ is then not identified --- only its ratio to the baseline is. $\beta$, $\omega$, the tax rate and the savings rate are common to the two variants; what differs is anything defined through $\eta$ or $X$.
\end{tablenotes}
\end{threeparttable}
\end{table}
